\documentclass{standalone}
\usepackage{circuitikz}
\usetikzlibrary{calc}
\definecolor{circuitnotemark}{HTML}{FE00FE}
\begin{document}
\begin{circuitikz}[american, line width=0.8pt]
\ctikzset{bipoles/length=1.2cm}
\foreach \x in {0,2,4,6,8,10,12} {\foreach \y in {0,-2,-4} {\fill[gray, opacity=0.35] (\x,\y) circle (1.1pt);}}
\node[gray, font=\scriptsize] at (0,0.84) {1};
\node[gray, font=\scriptsize] at (2,0.84) {2};
\node[gray, font=\scriptsize] at (4,0.84) {3};
\node[gray, font=\scriptsize] at (6,0.84) {4};
\node[gray, font=\scriptsize] at (8,0.84) {5};
\node[gray, font=\scriptsize] at (10,0.84) {6};
\node[gray, font=\scriptsize] at (12,0.84) {7};
\node[gray, font=\scriptsize, anchor=east] at (-0.84,0) {a};
\node[gray, font=\scriptsize, anchor=east] at (-0.84,-2) {b};
\node[gray, font=\scriptsize, anchor=east] at (-0.84,-4) {c};
\coordinate (a1) at (0,0);
\coordinate (c4) at (6,-4);
\coordinate (a7) at (12,0);
\draw (a1) node[ocirc]{} node[above left]{IN}; % line 2
\draw (a1) to[R, l_=$R_{1}$] (c4); % line 3
\draw (c4) to[R, l_=$R_{2}$] (a7); % line 4
\draw (a7) node[ocirc]{} node[above left]{OUT}; % line 5
\draw (a1) -- (a7); % line 7
\node[circ] at (a1) {};
\node[circ] at (a7) {};
\path (16,-4.29) rectangle (19.675,0.395); % line 9
\node[anchor=west, inner sep=0, circuitnotemark, font=\footnotesize] at (16,0) {X}; % line 9
\node[anchor=west, inner sep=0, circuitnotemark, font=\footnotesize] at (16,-0.325) {X}; % line 9
\node[anchor=west, inner sep=0, circuitnotemark, font=\footnotesize] at (16,-0.649) {X}; % line 9
\node[anchor=west, inner sep=0, circuitnotemark, font=\footnotesize] at (16,-0.974) {X}; % line 9
\node[anchor=west, inner sep=0, circuitnotemark, font=\footnotesize] at (16,-1.298) {X}; % line 9
\node[anchor=west, inner sep=0, circuitnotemark, font=\footnotesize] at (16,-1.623) {X}; % line 9
\node[anchor=west, inner sep=0, circuitnotemark, font=\footnotesize] at (16,-1.947) {X}; % line 9
\node[anchor=west, inner sep=0, circuitnotemark, font=\footnotesize] at (16,-2.272) {X}; % line 9
\node[anchor=west, inner sep=0, circuitnotemark, font=\footnotesize] at (16,-2.596) {X}; % line 9
\node[anchor=west, inner sep=0, circuitnotemark, font=\footnotesize] at (16,-2.921) {X}; % line 9
\node[anchor=west, inner sep=0, circuitnotemark, font=\footnotesize] at (16,-3.246) {X}; % line 9
\node[anchor=west, inner sep=0, circuitnotemark, font=\footnotesize] at (16,-3.57) {X}; % line 9
\node[anchor=west, inner sep=0, circuitnotemark, font=\footnotesize] at (16,-3.895) {X}; % line 9
\end{circuitikz}
\end{document}
